\documentclass[10pt,a4paper]{amsart}
\usepackage[margin=19mm]{geometry}
\usepackage{amsmath,amssymb,mathtools}
\usepackage[scr=rsfs,cal=euler]{mathalfa}
\usepackage{xfrac}
\usepackage[unicode,hidelinks,bookmarks=false]{hyperref}
\newcommand{\ie}{i.e.\ }
\newcommand{\cf}{cf.\ }
\newcommand{\resp}{respectively\ }
\newcommand{\Subsection}{\subsection}
\providecommand{\nobreakdash}{\nobreak\hbox{-}}
\setlength{\emergencystretch}{2em}
\multlinegap0pt
% math
\DeclareMathOperator{\sinc}{sinc}
\newcommand{\linf}[1][n\to\infty]{\liminf_{#1}}
\newcommand{\lsup}[1][n\to\infty]{\limsup_{#1}}
\newcommand{\dd}{\,\mathrm{d}}
\newcommand{\ddd}{\mathrm{d}}
\newcommand{\e}{\mathrm{e}}
\newcommand{\ccc}{\mathrm{c}}
\newcommand{\Tr}{\mathrm{Tr}}
\newcommand{\BM}{\mathrm{BM}}
\providecommand{\MR}{}\renewcommand{\MR}{\mathrm{MR}}
\newcommand{\UH}{\mathscr{L}_2(U,H)}
\newcommand{\one}{\mathbbm{1}}
\newcommand{\<}{\langle}
\newcommand{\?}{\rangle}
\newcommand{\pf}{p_{\alpha}}
\newcommand{\nn}{|\!|\!|}
\newcommand{\bigcdot}{\boldsymbol{\cdot}}
\newcommand{\der}[3][d]{\frac{\mathrm{#1}#2}{\mathrm{#1}#3}}
\newcommand{\into}{\hookrightarrow}
\newcommand{\calE}{\mathcal{E}}
\newcommand{\law}{\mathrm{Law}}
\DeclareMathOperator*{\esssup}{\textup{ess\,sup}}
\DeclareMathOperator*{\essinf}{\textup{ess\,inf}}
\newcommand{\supp}{\mathrm{supp}}
\newcommand{\loc}{\mathrm{loc}}
\newcommand{\ceqq}{\coloneqq}
\newcommand{\eqqc}{\eqqcolon}
\newcommand{\col}{\colon}
\newcommand{\indep}{\perp\!\!\!\perp} 
\renewcommand{\H}{\mathbb{H}} 
\renewcommand{\O}{\mathcal{O}} 
\newcommand{\bphieps}{\Psi^{\eps}} 
\newcommand{\wtveps}{%
\mspace{2mu}%
  \widetilde{\mspace{-2mu}\rule{0pt}{1.3ex}\smash[t]{v^\eps}}%
}
\newcommand{\y}{\Xi}
\newcommand{\gtn}{G_{\theta,n}}
\newcommand{\gvn}{G_{v,n}}
\newcommand{\pr}{\mathbb{P}}
\newcommand{\qr}{\mathbb{Q}}
\newcommand{\ft}{F_{\theta}}
\newcommand{\fvt}{f}
\newcommand{\gvtn}{g_n}
\newcommand{\fts}{F^{S}_{\theta}}
\newcommand{\fvs}{F^{S}_{v}}
\newcommand{\fv}{F_{v}}
\newcommand{\gt}{G_{\theta}}
\newcommand{\gv}{G_{v}}
\newcommand{\Progress}{\mathscr{P}}
\renewcommand{\l}{\langle}
\renewcommand{\r}{\rangle}
\newcommand{\pz}{\partial_3}
\newcommand{\pzz}{\partial_{33}} 
\newcommand{\DeltawR}{\Delta_{{\rm R}}^{{\rm w}}}

%spaces with boundary conditions
\newcommand{\BsD}{\prescript{}{D}{\Bs}}
\newcommand{\Ls}{\mathbb{L}}
\newcommand{\HsD}{\prescript{}{D}{\Hs}}
\newcommand{\Hps}{\prescript{}{ps}{H}}
\newcommand{\Hs}{\mathbb{H}}
\newcommand{\Hsd}{\prescript{}{D}{\mathbb{H}}}
\newcommand{\HD}{\prescript{}{D}{H}}
\newcommand{\Dd}{\prescript{}{D}{\Delta}}
\newcommand{\WD}{\prescript{}{D}{W}}
\newcommand{\BD}{\prescript{}{D}{B}}
\newcommand{\Bs}{\mathbb{B}}
\newcommand{\Xmhd}{\prescript{\Mhd}{D}{H}}
\newcommand{\Mhd}{\mathcal{M}}
\newcommand{\Tor}{\mathbb{T}}
\newcommand{\Borel}{\mathscr{B}}

\newcommand{\Br}{\mathrm{B}}
\newcommand{\h}{{\rm H}}
\newcommand{\n}{{\rm N}}
\newcommand{\Dr}{\Delta_{{\rm R}}^{{\rm w}}}
\newcommand{\Gforce}{\mathcal{G}}
\newcommand{\Force}{\mathcal{F}}
\newcommand{\forcetwo}{\mathcal{E}}
\newcommand{\force}{\mathcal{F}}
\renewcommand{\div}{\mathrm{div}}
\newcommand{\w}{w}
\newcommand{\Lb}{L^{2/(2\tilde{\beta}-1)}}
\renewcommand{\Tor}{\mathbb{T}^2}

% number systems
\def\N{{\mathbb N}}
\def\Z{{\mathbb Z}}
\def\Q{{\mathbb Q}}
\def\R{{\mathbb R}}
\def\C{{\mathbb C}} 

% probability notations
\newcommand{\E}{\mathsf{E}}
\newcommand{\B}{\mathbb{B}}
\newcommand{\W}{\mathbb{W}}
\newcommand{\BB}{\mathcal{B}}
\newcommand{\F}{\mathcal{F}}
\newcommand{\G}{\mathcal{G}}
\renewcommand{\P}{\mathbb{P}}
\newcommand{\hhp}{\mathrm{P}}
\newcommand{\PP}{\mathcal{P}}
\newcommand{\A}{\mathcal{A}}
\newcommand{\Filtr}{\mathbb{F}}
\newcommand{\X}{X_{\scriptscriptstyle\frac{1}{2}}}

% greek letters
\newcommand{\g}{\gamma}
\renewcommand{\d}{\delta}
\newcommand{\om}{\omega}
\newcommand{\Om}{\Omega}
\newcommand{\eps}{\varepsilon}

\usepackage{enumitem}
\theoremstyle{plain}
\newtheorem{theorem}{Theorem}
\newtheorem{lemma}[theorem]{Lemma}
\newtheorem{assumption}[theorem]{Assumption}
\title[Large deviations for stochastic evolution equations beyond the coercive case]{Large deviations for stochastic evolution equations beyond the coercive case: the critical-nonlinearity estimates for $\hat{F}$ (Lemma 3.5)}
\author{Esmée Theewis}
\date{Source: arXiv:2512.19501v2 (CC BY 4.0)}
\begin{document}
\maketitle

\noindent\emph{Source and licence.} Large deviations for stochastic evolution equations beyond the coercive case. Version arXiv:2512.19501v2 (CC BY 4.0), distributed under the Creative Commons Attribution 4.0 International licence, http://creativecommons.org/licenses/by/4.0/, as recorded in the arXiv licence metadata for this exact version and shown on the abstract page https://arxiv.org/abs/2512.19501v2. Full licence notice: licenses/arxiv-2512.19501-CC-BY-4.0.txt.\par\smallskip

\noindent\textit{Editorial scope.} Counted unit: the article's Lemma 3.5 (source0.tex lines 697--718, source label \texttt{lem: 3.8 extra F}), the lemma that supplies the estimates for the article's new nonlinearity $\hat{F}=\sum_i\hat{F}_i$ under the $\alpha$-(sub)criticality condition, delivered together with the article's complete original proof of it (source0.tex lines 720--762, one \texttt{proof} environment of 43 lines, adjacent to the statement, with no gap and no deferral; the article's printed numbering makes this Lemma 3.5 and its printed reference is to Lemma 3.5 in the published text). No mathematics is written, repaired or re-derived: statement and proof are the article's own text line for line, in the article's own notation ($\hat{\beta}_i$, $\hat{\rho}_i$, $\alpha_i$, $q$, $q'$, $\MR(0,T)$, $V_{\hat{\beta}}$, $V_\alpha$, $\<u,v\?$), and no line of the counted statement or of the counted proof is omitted, substituted or re-flowed. Required context retained uncounted: the article's Assumption 2.1, its extended critical variational setting (source0.tex lines 295--340, printed as Assumption 2.1, the assumption that the counted proof invokes in every one of its three parts), and the article's Lemma 3.4, the critical interpolation estimate (source0.tex lines 687--693, printed as Lemma 3.4), which is the one further estimate the counted proof invokes. Editorial formatting only, in addition: an AMS \texttt{amsart} preamble that restates the article's own macro definitions (including \texttt{\textbackslash MR}, \texttt{\textbackslash UH}, \texttt{\textbackslash nn}, \texttt{\textbackslash ceqq}, \texttt{\textbackslash dd}, \texttt{\textbackslash loc} and the article's number systems) and the theorem environments the retained lines use, with \texttt{enumitem} loaded for the article's own list formatting; sequential numbering of the retained environments, so that the retained assumption and the two retained lemmas are numbered 1, 2 and 3 in order of appearance, whereas the article prints them as Assumption 2.1, Lemma 3.4 and Lemma 3.5; the article's own section-relative equation numbering (its \texttt{\textbackslash numberwithin}) is not reproduced, so equations displayed inside the retained spans are numbered consecutively instead, and every cross-reference inside the retained text resolves to the wrapper's numbering. The article's preamble packages are otherwise the pinned AMS runtime; the article's own biblatex resource file, its class file if any, and its \texttt{.bbl} are neither shipped nor input. No citation occurs inside any retained line, so the wrapper carries no bibliography and no bibliography entry; no asset-inclusion command occurs inside any retained line, so the wrapper carries no asset.

\section*{Retained context: the extended critical variational setting}

\begin{assumption}[Local well-posedness conditions -- extended critical variational setting]\noindent\phantomsection\label{ass:critvarsettinglocal} 

\noindent
$A$ and $B$ are defined by 
\begin{equation}\label{eq:defAB}
A(t,v)=A_0(t,v)v-\hat{F}(t,v)-\hat{f}, \quad B(t,v)=B_0(t,v)v+G(t,v)+g, \qquad t\in\R_+,v\in V,
\end{equation}
and the following holds: 
\begin{enumerate}[label=\textit{(\arabic*)},ref=\textit{(\arabic*)}]
\item\label{it:AB mble}  $A_0\colon \R_+ \times H\to\mathscr{L}(V,V^*)$, $B_0\colon \R_+ \times H \to \mathscr{L}(V,\UH)) 
$  and $G\colon \R_+\times V\to\UH$ are Borel measurable and  $g\in L^2_\loc(\R_+;\UH)$. 

 \item\label{it:coercivelinear}  
  For all $n,T>0$, there exist  $\theta_{n,T},M_{n,T}>0$ such that for all $t\in[0,T],u\in H,v\in V$ with $\|u\|_H\leq n$:
  \begin{flalign*}
  &\qquad\qquad \<A_0(t,u)v,v\?-\tfrac{1}{2}\nn B_0(t,u)v\nn_H^2\geq \theta_{n,T}\|v\|_V^2-M_{n,T}\|v\|_H^2. \qquad &\text{\emph{(coercivity of the linear part)}}&
  \end{flalign*}
  
\item \label{it:growth AB} There exist an $m\in\N$ and $\rho_i\geq 0$, $\beta_i\in(\frac{1}{2},1)$ for $i\in\{1,\ldots, m\}$,  such that 
\begin{flalign}\label{eq: subcrit cond}
&\qquad\qquad (1+\rho_i)(2\beta_i-1)\leq 1,& \text{\emph{((sub)criticality)}}&
\end{flalign}
and for all  $n,T>0$ there exists a constant $C_{n,T}$ such that for all $t\in[0,T]$ and $u,v,w\in V$ with $\|u\|_H,\|v\|_H\leq n$, we have
\begin{align*}
\|A_0(t,u)w\|_{V^*}&\leq C_{n,T} \|w\|_V,\\
\|A_0(t,u)w-A_0(t,v)w\|_{V^*}&\leq C_{n,T}\|u-v\|_H\|w\|_V,\\
\nn B_0(t,u)w\nn_{H}&\leq C_{n,T} \|w\|_V,\\
\nn B_0(t,u)w-B_0(t,v)w\nn_{H}&\leq C_{n,T}\|u-v\|_H\|w\|_V,\\ 
\nn G(t,u)\nn_{H}&\leq C_{n,T}\textstyle{\sum_{i=1}^{m}} (1+\|u\|_{\beta_i}^{\rho_i+1}),\\
\nn G(t,u)-G(t,v)\nn_{H}&\leq C_{n,T}\textstyle{\sum_{i=1}^{m}} (1+\|u\|_{\beta_i}^{\rho_i}+\|v\|_{\beta_i}^{\rho_i})\|u-v\|_{\beta_i}.
\end{align*}
\item\label{it:hat F}  
$\hat{F}=\sum_{i=1}^{m}\hat{F}_i$ with $\hat{F}_i\colon\R_+\times V\to V_{\alpha_i}$ Borel measurable and $\alpha_i\in[0,\frac12]$,   
and there exist $\hat{\rho}_i>0$, $\hat{\beta}_i\in(\frac12,1]$ for $i\in\{1,\ldots, m\}$, such that
\begin{flalign}\label{eq: subcrit cond alpha}
&\qquad\qquad
  (1+\hat{\rho}_i)(2\hat{\beta}_i-1)\leq 1+2\alpha_i, & \text{\emph{($\alpha$-(sub)criticality)}}&
\end{flalign}
and for all $n,T>0$, $t\in[0,T]$ and $u,v\in V$ with $\|u\|_H,\|v\|_H\leq n$, 
\begin{align*} 
\|\hat{F}_i(t,u)\|_{\alpha_i}&\leq C_{n,T}  (1+\|u\|_{\hat{\beta}_i}^{\hat{\rho}_i+1}),\\
\|\hat{F}_i(t,u)-\hat{F}_i(t,v)\|_{\alpha_i}&\leq C_{n,T}(1+\|u\|_{\hat{\beta}_i}^{\hat{\rho}_i}+\|v\|_{\hat{\beta}_i}^{\hat{\rho}_i})\|u-v\|_{\hat{\beta}_i}.
\end{align*}
\item $\hat{f}=\sum_{i=1}^m \hat{f}_i$ with $\hat{f}_i\in L^{2/(1+2\alpha_i)}_{\loc}(\R_+;V_{\alpha_i})$. 
\end{enumerate}
\end{assumption}


\section*{Retained context: the critical interpolation estimate}

\begin{lemma}[Critical interpolation estimate]\label{lem:crit interpol}
Let  $\hat{\beta}\in (1/2,1]$ and $T>0$.    
Then for all $u\in \MR(0,T)$, it holds that 
\begin{equation*} 
  \|u\|_{L^{{2}/({2\hat{\beta}-1})}(0,T;V_{\hat{\beta}})}\leq   \|u\|_{L^{\infty}(0,T;H)}^{2-2\hat{\beta}}\|u\|_{L^{2}(0,T;V)}^{2\hat{\beta}-1}\leq \|u\|_{\MR(0,T)}.
\end{equation*}
\end{lemma} 


\section{The counted unit: Lemma 3.5 and its complete original proof}

\begin{lemma}\label{lem: 3.8 extra F}
 Let   $\hat{F}=\sum_{i=1}^m \hat{F}_i$ and $m,\alpha_i,\hat{\beta}_i,\hat{\rho}_i$ be as in Assumption \ref{ass:critvarsettinglocal}\ref{it:hat F}.  
  Let $n\in\R_+$, $T>0$ and  $\sigma>0$. 
  There exist  constants $M_T,\hat{C}_{n,T},\tilde{C}_{n,T},C_{n,T,\sigma}$ non-decreasing in $T$ and $n$, and constants $\epsilon_i>0$,  such that
  \begin{enumerate}[label=\emph{(\roman*)},ref=\textup{(\roman*)}]
  \item \label{it:emb4} for all $u\in\MR(0,T)$ with $\|u\|_{C([0,T];H)}\leq n$ and $i\in\{1,\ldots,m\}$:
  \begin{equation*}\label{eq: estimate hat F}
  \|\hat{F}_i(u)\|_{L^{2/(1+2\alpha_i)}(0,T;V_{\alpha_i})}  
   \leq \hat{C}_{n,T}(1+\|u\|_{L^2(0,T;V)}^{1+2\alpha_i}), 
  \end{equation*}
  \item \label{it:emb5} for all $u,v\in\MR(0,T)$ with $\|u\|_{C([0,T];H)},\|v\|_{C([0,T];H)}\leq n$ and $i\in\{1,\ldots,m\}$:
  \begin{align*}\label{eq: estimate hat F difference} 
  &\|\hat{F}_i(u)-\hat{F}_i(v)\|_{L^{2/(1+2\alpha_i)}(0,T;V_{\alpha_i})} \\  
  &\qquad\leq \tilde{C}_{n,T} \Big(T^{\epsilon_i}+\|u\|_{L^{{2}/({2\hat{\beta}_i-1})}(0,T;V_{\hat{\beta}_i})}^{\hat{\rho}_i}+\|v\|_{L^{{2}/({2\hat{\beta}_i-1})}(0,T;V_{\hat{\beta}_i})}^{ \hat{\rho}_i}\Big)\|u-v\|_{L^{{2}/({2\hat{\beta}_i-1})}(0,T;V_{\hat{\beta}_i})}, 
  \end{align*}
  \item \label{it:emb6} for all $u,v\in\MR(0,T)$ with $\|u\|_{C([0,T];H)},\|v\|_{C([0,T];H)}\leq n$:
  \begin{align*} 
  \Big|\int_0^T\<\hat{F}(u)-\hat{F}(v),u-v\?\dd t   \Big|
  &\leq \hat{C}_{\sigma,n,T} \int_0^T  (1+\|u\|_V^2+ \|v\|_V^2)\|u-v\|_H^2\dd t+\sigma\|u-v\|_{L^2(0,T;V)}^2. 
  \end{align*} 
  \end{enumerate}
\end{lemma}

\begin{proof}
Note that the constants $C_{n,T}$ in Assumption \ref{ass:critvarsettinglocal}\ref{it:hat F} can be taken non-decreasing in $T$ and $n$ without loss of generality. To omit subscripts, we provide the proof for $m=1$ and write $\alpha\ceqq \alpha_1$, $\hat{\beta}\ceqq\hat{\beta}_1$ and $\hat{\rho}\ceqq\hat{\rho}_1$. Applying the proof below for each $i$ and taking maxima or sums of the corresponding constants yields the case $m>1$. 
 
\ref{it:emb4}: Using subsequently Assumption \ref{ass:critvarsettinglocal}\ref{it:hat F}, H\"older's inequality, Lemma \ref{lem:crit interpol} and Young's inequality, we have
\begin{align*}
  \|\hat{F}(u)\|_{L^{2/(1+2\alpha)}(0,T;V_{\alpha})} 
  &\leq  C_{n,T} \big\|1+\|u\|_{\hat{\beta}}^{1+\hat{\rho}} \big\|_{L^{2/(1+2\alpha)}(0,T)} \\
   &=  {C}_{n,T}\big(T^{(1+2\alpha)/2}+\|u\|_{L^{2(1+\hat{\rho})/(1+2\alpha)}(0,T;V_{\hat{\beta}})}^{1+\hat{\rho}} \big)\\
   &\leq  \tilde{C}_{n,T}\big(1+\|u\|_{L^{2/(2\hat{\beta}-1)}(0,T;V_{\hat{\beta}})}^{1+\hat{\rho}} \big)\\
&\leq  \tilde{C}_{n,T}\big(1+\|u\|_{C([0,T];H)}^{(1+\hat{\rho})(2-2\hat{\beta})} \|u\|_{L^2(0,T;V)}^{(1+\hat{\rho})(2\hat{\beta}-1)} \big)\\
&\leq  \hat{C}_{n,T}\big(1+  \|u\|_{L^2(0,T;V)}^{1+2\alpha} \big),   
\end{align*}
where all constants are non-decreasing in $T$ and $n$. 
 
\ref{it:emb5}: 
Let $q=2(\hat{\rho}+1)/(1+2\alpha)$, $q' =2(\hat{\rho}+1)/((1+2\alpha)\hat{\rho})$ and note that by the (sub)criticality condition, we have $q=\hat{\rho}q'\leq {2}/({2\hat{\beta}-1})$. Combined with Assumption \ref{ass:critvarsettinglocal}\ref{it:hat F}, H\"older's inequality, Lemma \ref{lem:crit interpol} and Young's inequality (using that $2\hat{\beta}-1\in(0,1]$), this gives
\begin{align*}
\|&\hat{F}(u)-\hat{F}(v)\|_{L^{2/(1+2\alpha)}(0,T;V_{\alpha})}\leq C_{n,T} \big\|\big(1+\|u\|_{\hat{\beta}}^{\hat{\rho}}+\|v\|_{\hat{\beta}}^{ \hat{\rho}}\big)\|u-v\|_{\hat{\beta}} \big\|_{L^{2/(1+2\alpha)}(0,T)} \\
&\leq {C}_{n,T} \Big(T^{1/q'}+\|u\|_{L^{\hat{\rho}q'}(0,T;V_{\hat{\beta}})}^{\hat{\rho}}+\|v\|_{L^{\hat{\rho}q'}(0,T;V_{\hat{\beta}})}^{ \hat{\rho}}\Big)\|u-v\|_{L^{q}(0,T;V_{\hat{\beta}})}\\
&\leq \tilde{C}_{n,T} \Big(T^{1/q'}+\|u\|_{L^{{2}/({2\hat{\beta}-1})}(0,T;V_{\hat{\beta}})}^{\hat{\rho}}+\|v\|_{L^{{2}/({2\hat{\beta}-1})}(0,T;V_{\hat{\beta}})}^{ \hat{\rho}}\Big)\|u-v\|_{L^{{2}/({2\hat{\beta}-1})}(0,T;V_{\hat{\beta}})} 
\end{align*} 
where $C_{n,T}$ and $\tilde{C}_{n,T}$ can be chosen non-decreasing in $T$ and $n$. 

\ref{it:emb6}:  
We use the following complex interpolation estimates for $x\in V$, $y\in V_\alpha$ and $\theta\in[1/2,1]$:  
\begin{equation}\label{eq:interpol est alpha pair}
\begin{split}
&\|x\|_{\theta}\leq \|x\|_H^{2-2\theta}\|x\|_V^{2\theta-1}, \\   
&\<y,x\? \leq   \|y\|_{\alpha}\|x\|_{H}^{2\alpha}\|x\|_V^{1-2\alpha}, 
\end{split}
\end{equation}
where we use that $V_\theta=[H,V]_{2\theta-1}$ for $\theta\in[1/2,1]$. 
The   estimates of \eqref{eq:interpol est alpha pair}, Assumption \ref{ass:critvarsettinglocal}\ref{it:hat F} and Young's inequality  with $q= (1+\alpha-\hat{\beta})^{-1}$, $q'= (-\alpha+\hat{\beta})^{-1}$ yield
  \begin{align*} 
   \Big|\int_0^T\<\hat{F}(u)-\hat{F}(v),u-v\?\dd t    \Big|
  &\leq   \int_0^T \|\hat{F}(u)-\hat{F}(v)\|_\alpha\|u-v\|_H^{2\alpha}\|u-v\|_V^{1-2\alpha} \dd t \\
  &\leq  {C}_{n,T} \int_0^T  (1+\|u\|_{\hat{\beta}}^{\hat{\rho}}+ \|v\|_{\hat{\beta}}^{\hat{\rho}})\|u-v\|_H^{2\alpha+2-2\hat{\beta}}\|u-v\|_V^{-2\alpha+2\hat{\beta}} \dd t \\
  &\leq  \tilde{C}_{\sigma,n,T} \int_0^T  (1+\|u\|_{\hat{\beta}}^{\hat{\rho}q}+ \|v\|_{\hat{\beta}}^{\hat{\rho}q}) \|u-v\|_H^{2} \dd t +\sigma\|u-v\|_{L^2(0,T;V)}^2 \\
   &\leq \hat{C}_{\sigma,n,T} \int_0^T  (1+\|u\|_V^2+ \|v\|_V^2)\|u-v\|_H^2\dd t+\sigma\|u-v\|_{L^2(0,T;V)}^2. 
  \end{align*} 
 In the last line we  combine \eqref{eq:interpol est alpha pair} with the fact that $2\hat{\beta}-1\in(0,1]$ and 
$\hat{\rho}q\leq  {2}/({2\hat{\beta}-1})$ by virtue of  \eqref{eq: subcrit cond alpha}, and we use that    $\|u\|_{C([0,T];H)}\leq n$. Also, we note that $q\in(1,\infty)$, which is ensured by  \eqref{eq: subcrit cond alpha} and the fact that $\hat{\rho}>0$, $\alpha\geq 0$ and $\hat{\beta}\leq 1$. 
\end{proof}

\end{document}
